\chapter{Eksperymenty}
\label{cha:chapter003}

\section{Cel eksperymentów}

Eksperymenty miały na celu porównanie trzech metod klasyfikacji pod względem
dokładności na tych samych danych testowych oraz jakości predykcji w warunkach
rzeczywistych (inferencja z kamery).

\section{Podział danych}

Dane treningowe z pliku \texttt{sign\_mnist\_train.csv} zostały podzielone
stratyfikowanie na podzbiory:

\begin{itemize}
    \item \textbf{Treningowy (80\%)} --- używany do dopasowania parametrów modelu.
    \item \textbf{Walidacyjny (10\%)} --- opcjonalnie do strojenia hiperparametrów.
    \item \textbf{Testowy (10\%)} --- do pomiaru końcowej dokładności.
\end{itemize}

Plik \texttt{sign\_mnist\_test.csv} (7\,172 próbki) pełnił rolę niezależnego
zbioru testowego wymaganego przez dokumentację projektu.

\section{Procedura trenowania i ewaluacji}

Dla każdego klasyfikatora procedura była identyczna:

\begin{enumerate}
    \item Wczytanie \texttt{sign\_mnist\_train.csv} i normalizacja pikseli.
    \item Stratyfikowany podział 80/10/10.
    \item Trenowanie modelu na zbiorze treningowym.
    \item Zapis artefaktu do katalogu \texttt{artifacts/}.
    \item Ewaluacja na niezależnym zbiorze testowym.
\end{enumerate}

\subsection{Trenowanie SVC}

\begin{lstlisting}[style=styleBash, caption={Polecenie trenowania klasyfikatora SVC}, label=lst:train-svc]
uv run python main.py train \
    --classifier svc \
    --model artifacts/svc_model.pkl \
    --dataset data/raw/sign_mnist_train.csv
\end{lstlisting}

\subsection{Trenowanie Lasu Losowego}

\begin{lstlisting}[style=styleBash, caption={Polecenie trenowania Lasu Losowego}, label=lst:train-rf]
uv run python main.py train \
    --classifier rf \
    --model artifacts/rf_model.pkl \
    --dataset data/raw/sign_mnist_train.csv
\end{lstlisting}

\subsection{Trenowanie CNN}

\begin{lstlisting}[style=styleBash, caption={Polecenie trenowania CNN}, label=lst:train-cnn]
uv run python main.py train \
    --classifier cnn \
    --model artifacts/cnn_model.keras \
    --dataset data/raw/sign_mnist_train.csv
\end{lstlisting}

\subsection{Ewaluacja modeli}

\begin{lstlisting}[style=styleBash, caption={Ewaluacja modelu SVC na zbiorze testowym}, label=lst:eval]
uv run python main.py eval \
    --model artifacts/svc_model.pkl \
    --dataset data/raw/sign_mnist_test.csv
\end{lstlisting}

Wynikiem jest dokładność globalna oraz raport \texttt{classification\_report}
ze scikit-learn zawierający precyzję, czułość i miarę F1 dla każdej klasy.

\section{Test w warunkach rzeczywistych}

Po wytrenowaniu modeli przeprowadzono testy jakościowe z użyciem kamery
internetowej, zachowując warunki zestawione w tabeli~\ref{tab:warunki}.
Każdą z 24 liter testowano przez minimum 5 sekund, obserwując stabilność
predykcji w oknie OpenCV.

\begin{lstlisting}[style=styleBash, caption={Uruchomienie inferencji z kamery}, label=lst:infer-cam]
uv run python main.py infer \
    --source camera \
    --camera 0 \
    --model artifacts/svc_model.pkl
\end{lstlisting}

\subsection{Obserwacje jakościowe z inferencji kamerowej}

\begin{itemize}
    \item \textbf{SVC} --- predykcja stabilna w dobrym oświetleniu, wrażliwa
          na obrót dłoni i niejednolite tło.
    \item \textbf{Las Losowy} --- podobna stabilność jak SVC; tendencja do
          nadmiernej pewności siebie przy niskiej jakości obrazu.
    \item \textbf{CNN} --- najbardziej odporna predykcja; lepiej radzi sobie
          z lekkimi zmianami pozycji dłoni i oświetlenia.
\end{itemize}

We wszystkich modelach litery wizualnie podobne (np. M/N, R/U) powodowały
częstsze błędy klasyfikacji.